\documentclass[12pt, answers]{exam}
\usepackage{tikz}
\usepackage{amsmath}
\usepackage{amsfonts}
\usepackage{amssymb}
\usepackage{graphicx}
\usepackage{tabularray}
\usepackage{multirow}
\usepackage{xcolor}
\usepackage[most]{tcolorbox}
\tcbuselibrary{listings} 
\usepackage{xepersian}
\settextfont[Path = ../fonts/ ,Scale = 1.0 ,
BoldFont = *Bd ,
ItalicFont = *It ,
BoldItalicFont = *BdIt ,
Extension = .ttf
]{XB Yas}


\newtcblisting{codebox}[1]{
    enhanced,
    listing only,
    listing options={
        language=C,
        basicstyle=\ttfamily\small\latin,
        keywordstyle=\color{blue}\bfseries,
        commentstyle=\color{green!50!black},
        breaklines=true,
        numbers=left,
        numberstyle=\tiny\color{gray},
        tabsize=2
    },
    colback=white!95!gray,
    colframe=black!70,
    boxrule=0.5mm,
    arc=2mm,
    attach boxed title to bottom center={yshift=3mm}, 
    boxed title style={
        colback=black!10,
        colframe=black!70,
        boxrule=0.5mm,
        arc=1mm
    },
    title={\textbf{#1}},
    coltitle=black
}

\renewcommand{\solutiontitle}{\noindent\textbf{پاسخ:}\par}
\newcommand{\class}{اصول سیستم‌های عامل}
\newcommand{\term}{نیم‌سال اول ۰۴-۰۵}
\newcommand{\college}{دانشکده ریاضی و آمار}
\newcommand{\prof}{استاد: دکتر رفیعی}
\newcommand{\group}{\textcolor{red}{نام دانشجو - شماره دانشجویی}}


\begin{document}
\pointpoints{نمره}{نمره}

\pagestyle{head}
\firstpageheader{}{}{}
\runningheader{صفحه \thepage\ از \numpages}{}{\class}
\runningheadrule
\begin{tabular}{p{.7\textwidth} l}
    \multicolumn{2}{c}{\textbf{به نام خدا}}                                            \\
                    & \multirow{2}{*}{\includegraphics[scale=0.2]{../UILOGO.png}} \\
    \textbf{\class} &                                                                  \\
    \textbf{\term}  &                                                                  \\
    \textbf{\prof}  &                                                                  \\
    \textbf{\group} &                                                                  \\
                    & \textbf{\college}                                                \\
    % & \textbf{\todaydate} \\
    % & \textbf{\deadline} \\ \\
\end{tabular}\\
\rule[1ex]{\textwidth}{.1pt}

\begin{center}
    \textbf{تمرین سری چهارم}
\end{center}


\begin{questions}
    \question
    \begin{parts}
        \part[3] توزیع‌کننده یکی از مولفه‌های \dots \
        است. این مولفه در واقع یک \dots \
        است که کنترل \dots \ را در اختیار فرآیند انتخاب شده توسط زمان‌بند CPU
        قرار می‌دهد.

        \part[2] دو مورد از وظایف عملکردی توزیع‌کننده را نام ببرید.

        \part[3] آیا سرعت این مولفه بر روی عملکرد کلی سیستم‌عامل تاثیر می‌گذارد؟ چرا؟
    \end{parts}
    \begin{solution}
        \textcolor{red}{جواب خود را اینجا بنویسید.}
    \end{solution}

    \pagebreak
    \question
    \begin{parts}
        \part[1] چرا مدت زمان انفجارهای \lr{CPU} در سیستم اندازه‌گیری می‌شود؟
        \part[1] آیا صف فرآیندهای آماده (\lr{Ready Queue}) لزوماً باید به صورت \lr{FIFO} پیاده‌سازی شود؟
        \part[2] دو مورد از ساختارهای داده‌ای که می‌توان برای پیاده‌سازی صف آماده استفاده کرد را نام ببرید.

    \end{parts}

    \begin{solution}
        \textcolor{red}{جواب خود را اینجا بنویسید.}
    \end{solution}

    \pagebreak
    \question
    \begin{parts}
        \part[1] اگر زمان‌بندی فقط در شرایط «تغییر از اجرا به انتظار» و «خاتمه فرآیند» صورت گیرد، به آن طرح زمان‌بندی \dots \ می‌گوییم.
        \part[1] اکثر سیستم‌عامل‌های مدرن از چه مکانیزمی برای زمان‌بندی فرآیندهای کرنل استفاده می‌کنند؟
        \part[1] در اسلایدها خواندیم که بخش‌هایی از کد وجود دارند که
        نمی‌توانند توسط چند فرآیند به صورت همروند استفاده شوند. از این رو
        \dots \ ها هنگام ورود به این بخش غیرفعال و به هنگام خروج فعال می‌شوند.
    \end{parts}

    \begin{solution}
        \textcolor{red}{جواب خود را اینجا بنویسید.}
    \end{solution}


    \pagebreak
    \question در این سوال قصد داریم تا دو مشکل مربوط به
    \textbf{اجرای قبضه‌ای}
    فرآیندها را مورد بررسی قرار دهیم (به صفحه ۱۷ اسلاید ۱۳ مراجعه شود).
    \begin{parts}
        \part[10] (\textbf{مشکل ۱: از دست رفتن ورودی}) فرض کنید که دو فرایند داریم. فرایند ۱ دائم اطلاعات مربوط به \lr{I/O} را از یک دستگاه دریافت می‌کند و آن را در یک بافر مشترک قرار می‌دهد. فرایند ۲ نیز دائم داده‌ها را از این بافر می‌خواند و پردازش می‌کند.
        در صورت اجرای قبضه‌ای این دو فرایند، چه مشکلی ممکن است رخ دهد؟
        لطفاً با استفاده از کد زیر که نمایانگر هر یک از این دو فرایند است، مشکل را توضیح دهید.

        دقت داشته باشید که متغیر
        \lr{buffer} بین هر دو فرآیند مشترک است.

        متغیر \lr{new\_data} نیز داده جدیدی است که از دستگاه \lr{I/O} دریافت شده است.
        \begin{tcbraster}[raster columns=2, raster column skip=0.5cm, raster equal height]

            \begin{codebox}{فرآیند ۱}
while (1) {
    // Producer
    buffer = new_data;
}
            \end{codebox}

            \begin{codebox}{فرآیند ۲}
void kernel_process() {
    // Consumer
    int data = buffer;
    process(data);
}
            \end{codebox}

        \end{tcbraster}

        \vspace{1cm}
        \begin{solution}
            \textcolor{red}{جواب خود را اینجا بنویسید.}
        \end{solution}


        \pagebreak
        \part[10] (\textbf{مشکل ۲: بازنویسی خروجی}) فرض کنید که یک تابع به نام \lr{increment} داریم که یک شمارنده مشترک
        (counter)
        را افزایش می‌دهد. پیاده‌سازی این تابع را در زیر می‌بینید.
        این پیاده‌سازی تقریبا معادل به اتفاقاتی است که هنگام اجرای
        \lr{counter++}
        در زبان \lr{C} رخ می‌دهد.
        حالا فرض کنید این تابع دوبار توسط تابع
        main
        صدا زده شود. چه مشکلی ممکن است رخ دهد؟ توضیح دهید.
        \begin{tcbraster}[raster columns=2, raster column skip=0.5cm, raster equal height]

            \begin{codebox}{تابع افزایش مقدار شمارنده}
void increment() {
    int temp;
    temp = counter;
    temp = temp + 1;
    counter = temp;
}
            \end{codebox}

            \begin{codebox}{تابع main}
int main() {
    // Shared counter = 5
    increment();
    increment();
}
            \end{codebox}



        \end{tcbraster}
    \end{parts}

    \begin{solution}
        \textcolor{red}{جواب خود را اینجا بنویسید.}
    \end{solution}



    \pagebreak
    \question[20]
    سه فرآیند
    \lr{P3, P2, P1}
    با زمان اجرا و زمان ورود زیر را در نظر بگیرید و به موارد خواسته شده پاسخ دهید.
    (\lr{quantum = 2})

    \vspace{0.2cm}
    \begin{minipage}[c]{0.3\textwidth}
        \begin{enumerate}
            \item ﺑﻬﺮﻩ ﻭﺭﯼ CPU
            \item ﻣﯿﺎﻧﮕﯿﻦ ﺯﻣﺎﻥ ﺑﺮﮔﺸﺖ
            \item ﻣﯿﺎﻧﮕﯿﻦ ﺯﻣﺎﻥ ﺍﻧﺘﻈﺎﺭ
            \item ﻣﯿﺎﻧﮕﯿﻦ ﺯﻣﺎﻥ ﭘﺎﺳﺦ
        \end{enumerate}
    \end{minipage}
    \begin{minipage}[c]{0.7\textwidth}
        \centering
        \begin{tblr}{|c|c|c|}
            \hline
            زمان اجرا & زمان ورود & فرآیند \\ \hline
            $6$       & $0$       & $P1$   \\ \hline
            $3$       & $1$       & $P2$   \\ \hline
            $4$       & $6$       & $P3$   \\ \hline
            $2$       & $6$       & $P4$   \\
            \hline
        \end{tblr}
    \end{minipage}

    \begin{solution}
        \textcolor{red}{جواب خود را اینجا بنویسید.}
    \end{solution}

    \pagebreak
    \question
    \begin{parts}
        \part[5] تفاوت دیدگاه کاربر و دیدگاه سیستم در ارزیابی الگوریتم‌های زمان‌بندی چیست؟
        \part[5] منظور از زمان بازگشت \lr{(Turnaround Time)} چیست و چگونه محاسبه می‌شود؟
    \end{parts}

    \begin{solution}
        \textcolor{red}{جواب خود را اینجا بنویسید.}
    \end{solution}

    \pagebreak
    \question
    \begin{parts}
        \part[5] اگر در الگوریتم زمان‌بندی نوبتی-گردشی (RR)، \textbf{کوانتوم زمانی} بزرگ‌تر از زمان اجرای طولانی‌ترین فرآیند موجود در سیستم باشد، این الگوریتم به چه سیاستی تبدیل می‌شود؟


        \part[5] برای جلوگیری از پدیده \textbf{گرسنگی} در زمان‌بندی اولویت‌دار، از چه مکانیزمی استفاده می‌شود؟


        \part[5] چرا الگوریتم \textbf{SJF} از نظر تئوری یک الگوریتم \textbf{بهینه} محسوب می‌شود؟


        \part[5] کوچک انتخاب کردن بیش از حد \textbf{کوانتوم زمانی} در الگوریتم RR چه پیامدی برای سیستم دارد؟

    \end{parts}

    \begin{solution}
        \textcolor{red}{جواب خود را اینجا بنویسید.}
    \end{solution}

    \pagebreak
    \question
    \begin{parts}
        \part[10] با فرض استفاده از فرمول $S_n = \alpha T_{n-1} + (1-\alpha) S_{n-1}$ برای پیش‌بینی انفجار بعدی در SJF، اگر $\alpha = 0\.1$ تنظیم شود و یک فرآیند پس از مدت‌ها اجرای کوتاه، ناگهان یک انفجار بسیار طولانی داشته باشد، واکنش سیستم را تحلیل کنید.


        \part[10] تحلیل کنید که چرا در الگوریتم \textbf{MLFQ}، فرآیندهای I/O Bound در صفوف بالا باقی می‌مانند اما فرآیندهای CPU Bound به صفوف پایین رانده می‌شوند؟

    \end{parts}

    \begin{solution}
        \textcolor{red}{جواب خود را اینجا بنویسید.}
    \end{solution}
\end{questions}

\end{document}
